\documentstyle[twoside,12pt]{article}
\topmargin-0.5in
\textheight9.0in
\oddsidemargin0in
\evensidemargin0in
\textwidth6.5in
\pagestyle{myheadings}
\begin{document}
\input{psfig}
\input{rotate}
\newbox\rbox
% set up general text format for slides
{\raggedright
{\huge  
% end set up
\setcounter{page}{101}
\begin{center}
Section Outline
\end{center}
\vspace{0.5in}
\begin{enumerate}
\item Non-Constant Variation
\item Outliers
\end{enumerate}
\newpage
\markboth{Non-Constant Variation}{Non-Constant Variation}
\begin{center}
Non-Constant Variation
\end{center}
\vspace{0.75in}
Non-constant variation across the levels of the predictor
variables violates one of the usual assumptions of least
squares regression. \\
\vspace{0.5in}
However non-constant variation is often found in data
used for building regression models as an inherent part
of the measurement or other type of process. \\
\vspace{0.5in}
Because non-constant variation is frequently encountered, techniques 
to improve the validity of the standard assumptions and techniques 
based on alternative assumptions have been developed to lessen
the impact of this problem. 
\newpage
\begin{center}
Effect of Non-Constant Standard Deviation on Function Fitting
\end{center}
\vspace{0.75in}
If all of the other assumptions for the analysis are met,
except for constant standard deviation across all combinations
of the predictors, the usual parameter estimates will still
be correct on average, or unbiased. \\
\vspace{0.5in}
However, the deviation of the fitted function from the true function
will be larger, on average, than that from other techniques, unlike 
the case of constant standard deviation.
\newpage
\begin{center}
Effect of Non-Constant Standard Deviation on Function Fitting
\end{center}
\vspace{0.5in}
The least squares estimates are found by minimizing
the sum of the squared differences of the predicted values and the observed $y$'s, Q,
to find the $\hat{\beta}$'s.
\vspace{0.35in}
\[
Q=\sum_{i=1}^{n}[y_i-f(x_{1i},x_{2i},\ldots,x_{ki};\hat{\beta}_1,\hat{\beta}_2,\ldots,\hat{\beta}_p)]^2
\] \\
\vspace{0.5in}
All deviations count the same using the LSS fitting criterion, even if the variation
differs across the data. \\
\vspace{0.5in}
However, where the variation is larger, the deviations of the data from their corresponding
means will also be larger, indicating that the regression function should not be constrained
to be as near that data as it is to data that varies less.
\newpage
\begin{center}
Effect of Non-Constant Standard Deviation on Prediction
\end{center}
\vspace{0.75in}
In addition to increasing the variance of estimating the regression
function, non-constant standard deviation impacts the uncertainty of
predictions, if unaccounted for. \\
\vspace{0.5in}
Predictions made where variation is relatively low, will tend to have larger 
stated uncertainties than necessary. Predictions made where uncertainty is 
relatively large will have smaller stated uncertainties than necessary.
%\newpage
\begin{figure}
\psfig{file=hetsc1.ps,width=6.5in}
\end{figure}
\newpage
\begin{center}
Improved Fitting Methods
\end{center}
\vspace{0.75in}
There are basically two approaches to getting improved parameter
estimates for data with non-constant standard deviations:
\vspace{0.5in}
\begin{enumerate}
\item transformation of the data so it meets the standard assumptions, and
\vspace{0.15in}
\item use of weights in parameter estimation to account for the unequal 
standard deviations.
\end{enumerate}
%\newpage
\begin{figure}
\psfig{file=hetsc2.ps,width=6.5in}
\end{figure}
\newpage
\begin{center}
Transformations
\end{center}
\vspace{0.6in}
The basic steps for using transformations to handle data with
unequal subpopulation standard deviations are:
\vspace{0.4in}
\begin{enumerate}
\item Transform the response variable to equalize the variation
across the levels of the predictor variables
\vspace{0.15in}
\item Transform the predictor variables if necessary to attain
or restore a simple functional form for the regression function.
\vspace{0.15in}
\item Fit and validate the model in the transformed variables.
\vspace{0.15in}
\item Transform the function back into the original units, if 
necessary, using the inverse of the transformation originally 
applied to the response variable.
\end{enumerate}
\newpage
\begin{center}
Transformations
\end{center}
\vspace{0.75in}
Appropriate transformations to stabilize the 
variability may be suggested by scientific knowledge
or selected using the data. \\
\vspace{0.5in}
Three transformations that are often effective is
equalizing the standard deviations across the
values of the predictor are:
\vspace{0.5in}
\begin{enumerate}
\item $\sqrt{y}$
\item $\log(y)$, and
\item $1/y$.
\end{enumerate}
%\newpage
\begin{figure}
\psfig{file=hetsc3.ps,width=6.5in}
\end{figure}
%\newpage
\begin{figure}
\psfig{file=hetsc4.ps,width=6.5in}
\end{figure}
%\newpage
\begin{figure}
\psfig{file=hetsc5.ps,width=6.5in}
\end{figure}
%\newpage
\begin{figure}
\psfig{file=hetsc6.ps,width=6.5in}
\end{figure}
%\newpage
\begin{figure}
\psfig{file=hetsc7.ps,width=6.5in}
\end{figure}
%\newpage
\begin{figure}
\psfig{file=hetsc8.ps,width=6.5in}
\end{figure}
\newpage
\begin{center}
Weighted Least Squares
\end{center}
\vspace{0.5in}
Weighted least squares estimates are found by minimizing:
\vspace{0.5in}
\[
Q=\sum_{i=1}^{n}w_i[y_i-f(x_{1i},x_{2i},\ldots,x_{ki};\hat{\beta}_1,\hat{\beta}_2,\ldots,\hat{\beta}_p)]^2
\] \\
\vspace{0.25in}
where  \\
\vspace{0.25in}
\[ w_i \propto \frac{1}{\sigma^2_i} \] \\
\vspace{0.5in}
These relative weights give optimal results, when known.
\newpage
\begin{center}
Weighted Least Squares
\end{center}
\vspace{0.5in}
Unfortunately the true weights are rarely known, so they
have to be estimated. \\
\vspace{0.5in}
The obvious way to estimate the $w_i$, if there are replicates
in the data, is
\[ w_i = \frac{1}{\hat{\sigma}^2_i} = \frac{1}{s^2_i} = 
         [\frac{\sum_{i=1}^{n_i}(y_{ij}-\bar{y}_i)^2}{n_i-1}]^{-1} \]
for the $i^{th}$ set of replicates in the data set. \\
\vspace{0.5in}
However, this rarely works well because the weights are
extremely variable when estimate this way.
\newpage
\begin{center}
Weighted Least Squares
\end{center}
\vspace{0.5in}
A better strategy for estimating the weights is to find a function which relates
$s_i^2$ to $x_i$. \\
\vspace{0.5in}
If \[ s_i^2 \propto f(x_{1i},x_{2i},\ldots,x_{ki}) \] \\
then use \\
\[ w_i = \frac{1}{f(x_{1i},x_{2i},\ldots,x_{ki})} \] \\
as the weights. \\
\vspace{0.5in}
One model that {\em often} works well for modeling the variances is
\[ s_i^2 \propto x_i^c \]
%\newpage
\begin{figure}
\psfig{file=hetsc9.ps,width=6.5in}
\end{figure}
\newpage
\begin{center}
Estimation of the Weights - Replicates Available
\end{center}
\vspace{0.5in}
To estimate the weights using the power function
shown above, fit the function
\[ \log(s_i^2) = \beta_1 + \beta_2\log(x_i) \]
to the variances from each set of replicates in the
data. \\
\vspace{0.5in}
The use $\hat{c} = \hat{\beta}_2$ (the slope of the fit)
to estimate $c$, and \\
\[ w_i = \frac{1}{x^{\hat{c}}_i} \] \\
to as the weights. \\
\vspace{0.5in}
Check the residuals from the fit used to estimate $c$ just
to make sure everything looks reasonable.  The fit does not
have to meet the standards usually used, however. 
\newpage
\begin{center}
Estimation of the Weights - No Replicates Available
\end{center}
\vspace{0.5in}
If there are few or no replicates in the data, divide the
data into several ranges in which the responses have 
fairly similar means. \\
\vspace{0.5in}
Treat each range as replicates and
compute $\bar{x_i}$ and $s_i^2$ for each range. \\
\vspace{0.5in}
Then fit 
\[ \log(s_i^2) = \beta_1 + \beta_2\log(\bar{x}_i) \]
and define the weights by
\[ w_i = \frac{1}{x^{\hat{\beta}_2}_i} \]
%\newpage
\begin{figure}
\psfig{file=hetsc10.ps,width=6.5in}
\end{figure}
\clearpage
\newpage
\begin{center}
Output from Fit for Weight Estimation
\end{center}
\vspace{0.5in}
{\large
\begin{verbatim}

N = 10

Residual Standard Error = 0.7821

Multiple R-Square = 0.8046

F-statistic = 32.9428 on 1 and 8 df, p-value = 4e-04

             coef std.err  t.stat p.value 
Intercept  0.5554  1.3785  0.4029  0.6976
        X  2.0414  0.3557  5.7396  0.0004 
\end{verbatim}
}
%\clearpage
%\newpage
%\begin{center}
%Output from Fit for Weight Estimation
%\end{center}
%\vspace{0.5in}
%{\large
%\begin{verbatim}
%
%N = 10
%
%Residual Standard Error = 1.08
%
%Multiple R-Square = 0.806
%
%F-statistic = 33.2333 on 1 and 8 df, p-value = 4e-04
%
%             coef std.err  t.stat p.value 
%Intercept -3.2265  1.9035 -1.6950  0.1285
%        X  2.8312  0.4911  5.7648  0.0004
%\end{verbatim}
%}
\clearpage
\newpage
\begin{figure}
\psfig{file=hetsc11.ps,width=6.5in}
\end{figure}
%\newpage
\begin{figure}
\psfig{file=hetsc12.ps,width=6.5in}
\end{figure}
\clearpage
\newpage
\begin{center}
Weighted Residuals
\end{center}
\vspace{0.5in}
One complication with weighted analyses is the
fact that the distribution of the residuals can 
vary substantially with the different values of the 
predictor variables. \\
\vspace{0.5in}
This necessitates the use of weighted residuals when
plotting residuals. \\
\vspace{0.5in}
The weighted residuals are given by
\[ e_i = \sqrt{w_i}(y_i-f(x_{1i},\ldots,x_{ki};\hat{\beta}_1,\ldots,\hat{\beta}_p)) \]
\clearpage
%\newpage
\begin{figure}
\psfig{file=hetsc13.ps,width=6.5in}
\end{figure}
%\newpage
\begin{figure}
\psfig{file=hetsc14.ps,width=6.5in}
\end{figure}
%\newpage
\begin{figure}
\psfig{file=hetsc15.ps,width=6.5in}
\end{figure}
%\newpage
\begin{figure}
\psfig{file=berger1.ps,width=6.5in}
\end{figure}
%\newpage
\begin{figure}
\psfig{file=berger2.ps,width=6.5in}
\end{figure}
\clearpage
\newpage
\begin{center}
AK Pipeline Data \\
Output from  Unweighted Fit
\end{center}
\vspace{0.5in}
{\large
\begin{verbatim}

N = 107

Residual Standard Error = 6.080924

Multiple R-Square = 0.8941251

F-statistic = 886.7366 on 1 and 105 df, p-value = 0

               coef    std.err    t.stat      p.value 
Intercept 4.9936799 1.12565780  4.436233 2.263517e-05
        X 0.7311111 0.02455195 29.778123 0.000000e+00
\end{verbatim}
}
\clearpage
%\newpage
\begin{figure}
\psfig{file=berger3.ps,width=6.5in}
\end{figure}
%\newpage
\begin{figure}
\psfig{file=berger4.ps,width=6.5in}
\end{figure}
%\newpage
\begin{figure}
\psfig{file=berger5.ps,width=6.5in}
\end{figure}
%\newpage
\begin{figure}
\psfig{file=berger6.ps,width=6.5in}
\end{figure}
%\newpage
\begin{figure}
\psfig{file=berger7.ps,width=6.5in}
\end{figure}
\clearpage
\newpage
\begin{center}
AK Pipeline Data \\
Output from Fit on Transformed Variables
\end{center}
\vspace{0.5in}
{\large
\begin{verbatim}

N = 107

Residual Standard Error = 0.1682604

Multiple R-Square = 0.9337104

F-statistic = 1478.958 on 1 and 105 df, p-value = 0

               coef    std.err    t.stat      p.value 
Intercept 0.2813838 0.08092894  3.476924 0.0007390395
        X 0.8851754 0.02301714 38.457221 0.0000000000
\end{verbatim}
}
\clearpage
%\newpage
\begin{figure}
\psfig{file=berger8.ps,width=6.5in}
\end{figure}
%\newpage
\begin{figure}
\psfig{file=berger9.ps,width=6.5in}
\end{figure}
%\newpage
\begin{figure}
\psfig{file=berger10.ps,width=6.5in}
\end{figure}
\clearpage
\newpage
\begin{center}
AK Pipeline Data \\
Output from Weight Estimation Fit 
\end{center}
\vspace{0.5in}
{\large
\begin{verbatim}

N = 27

Residual Standard Error = 0.8545392

Multiple R-Square = 0.6286482

F-statistic = 42.32161 on 1 and 25 df, p-value = 8.178602e-07

               coef   std.err    t.stat      p.value 
Intercept -2.696496 0.8250090 -3.268444 3.140569e-03
        X  1.522101 0.2339712  6.505506 8.178602e-07
\end{verbatim}
}
\clearpage
%\newpage
\begin{figure}
\psfig{file=berger11.ps,width=6.5in}
\end{figure}
%\newpage
\begin{figure}
\psfig{file=berger12.ps,width=6.5in}
\end{figure}
%\newpage
\begin{figure}
\psfig{file=berger13.ps,width=6.5in}
\end{figure}
\clearpage
\newpage
\begin{center}
AK Pipeline Data \\
Output from Weighted Fit 
\end{center}
\vspace{0.5in}
{\large
\begin{verbatim}

N = 107

Residual Standard Error = 0.3646

Multiple R-Square = 0.9235

F-statistic = 1267.024 on 1 and 105 df, p-value = 0

            coef std.err  t.stat p.value 
Intercept 2.3523  0.5431  4.3312       0
        X 0.8064  0.0227 35.5953       0
\end{verbatim}
}
\clearpage
%\newpage
\begin{figure}
\psfig{file=berger14.ps,width=6.5in}
\end{figure}
%\newpage
\begin{figure}
\psfig{file=chwirut1.ps,width=6.5in}
\end{figure}
%\newpage
\begin{figure}
\psfig{file=chwirut2.ps,width=6.5in}
\end{figure}
\clearpage
\newpage
\begin{center}
Ultrasonic Calibration Data \\
Output from Unweighted Exp/Linear Fit
\end{center}
\vspace{0.5in}
{\large
\begin{verbatim}
Parameters:

        Value  Std. Error  t value 
b1 0.19027400 0.021938300  8.67312
b2 0.00613137 0.000345001 17.77200
b3 0.01053100 0.000792818 13.28300

Residual standard error: 3.36167 on 211 degrees of freedom

Correlation of Parameter Estimates:

       b1     b2 
b2  0.839       
b3 -0.950 -0.949
\end{verbatim}
}
\clearpage
%\newpage
\begin{figure}
\psfig{file=chwirut3.ps,width=6.5in}
\end{figure}
%\newpage
\begin{figure}
\psfig{file=chwirut4.ps,width=6.5in}
\end{figure}
%\newpage
\begin{figure}
\psfig{file=chwirut5.ps,width=6.5in}
\end{figure}
%\newpage
\begin{figure}
\psfig{file=chwirut6.ps,width=6.5in}
\end{figure}
%\newpage
\begin{figure}
\psfig{file=chwirut7.ps,width=6.5in}
\end{figure}
\clearpage
\newpage
\begin{center}
Ultrasonic Calibration Data \\
Output from Fit to Transformed Data 
\end{center}
\vspace{0.5in}
{\large
\begin{verbatim}
Parameters:

        Value Std. Error   t value 
b1 -0.0154274 0.00861101  -1.79159
b2  0.0806725 0.00150574  53.57670
b3  0.0638570 0.00288001  22.17250

Residual standard error: 0.29715 on 211 degrees of freedom

Correlation of Parameter Estimates:

       b1     b2 
b2  0.793       
b3 -0.960 -0.899
\end{verbatim}
}
\clearpage
%\newpage
\begin{figure}
\psfig{file=chwirut8.ps,width=6.5in}
\end{figure}
%\newpage
\begin{figure}
\psfig{file=chwirut9.ps,width=6.5in}
\end{figure}
\clearpage
\newpage
\begin{center}
Ultrasonic Calibration Data \\
Output from Weight Estimation Fit
\end{center}
\vspace{0.5in}
{\large
\begin{verbatim}

N = 22

Residual Standard Error = 0.6099457

Multiple R-Square = 0.6712342

F-statistic = 40.83357 on 1 and 20 df, p-value = 3.10602e-06

               coef   std.err     t.stat      p.value 
Intercept  2.536866 0.1919360  13.217253 2.421219e-11
        X -1.112763 0.1741382  -6.390115 3.106020e-06
\end{verbatim}
}
\clearpage
%\newpage
\begin{figure}
\psfig{file=chwirut10.ps,width=6.5in}
\end{figure}
%\newpage
\begin{figure}
\psfig{file=chwirut11.ps,width=6.5in}
\end{figure}
%\newpage
\begin{figure}
\psfig{file=chwirut12.ps,width=6.5in}
\end{figure}
\clearpage
\newpage
\begin{center}
Ultrasonic Calibration Data \\
Output from Weighted Fit
\end{center}
\vspace{0.5in}
{\large
\begin{verbatim}
LEAST SQUARES NON-LINEAR FIT
SAMPLE SIZE N =      214
MODEL--Y =EXP(-B1*X)/(B2+B3*X)
REPLICATION CASE
REPLICATION STANDARD DEVIATION =     0.3281762600D+01
REPLICATION DEGREES OF FREEDOM =         192
NUMBER OF DISTINCT SUBSETS     =          22
  
  
FINAL PARAMETER 
      ESTIMATES     (APPROX. ST. DEV.)    T VALUE
1  B1  0.143378       (0.1476E-01)          9.7
2  B2  0.518479E-02   (0.4191E-03)          12.
3  B3  0.125719E-01   (0.7508E-03)          17.
  
RESIDUAL    STANDARD DEVIATION =         4.2653684616
RESIDUAL    DEGREES OF FREEDOM =         211
REPLICATION STANDARD DEVIATION =         3.2817625999
REPLICATION DEGREES OF FREEDOM =         192
LACK OF FIT F RATIO =  8.6545 = THE 100.0000% POINT OF THE
F DISTRIBUTION WITH     19 AND    192 DEGREES OF FREEDOM
\end{verbatim}
}
\clearpage
%\newpage
\begin{figure}
\psfig{file=chwirut13.ps,width=6.5in}
\end{figure}
\clearpage
\newpage
\markboth{Outliers}{Outliers}
\begin{center}
Outliers
\end{center}
\vspace{0.5in}
In a broad sense, outliers are points that do 
not follow the general pattern or structure in the data. \\
\vspace{0.5in}
If present in the data, outliers can unduly affect
the model-building process and invalidate
predictions or calibrations. \\
\newpage
\begin{center}
More Specifically $\ldots$
\end{center}
\vspace{0.5in}
\begin{description}
\item[outlier] - an observed value of the response variable that
appears to be unusually large or small relative to its apparent population.
\vspace{0.15in}
\item[high leverage point] - a point from the observed set of predictor
variables that is unusual in size or structure relative to the other 
predictor variable points. 
\vspace{0.15in}
\item[influential observation] - an observation that has a large effect on
the model derived from the data.
\end{description}
\vspace{0.5in}
Note: Not all outliers or high leverage points are influential observations.
%\newpage
\begin{figure}
\psfig{file=infl1.ps,height=9.0in}
\end{figure}
\clearpage
\begin{center}
Regression Output for Data with Outlier
\end{center}
{\large
\begin{verbatim}

N = 10

Residual Standard Error = 2.6986

Multiple R-Square = 0.9252

F-statistic = 98.9308 on 1 and 8 df, p-value = 0

            coef std.err t.stat p.value 
Intercept 6.6162  1.8435 3.5889  0.0071
    x     2.9552  0.2971 9.9464  0.0000
\end{verbatim}
}
\vspace{0.5in}
\begin{center}
Regression Output for Data with Outlier Omitted
\end{center}
{\large
\begin{verbatim}

N = 9

Residual Standard Error = 1.4547

Multiple R-Square = 0.9803

F-statistic = 349.0913 on 1 and 7 df, p-value = 0

            coef std.err  t.stat p.value 
Intercept 5.6882  1.0146  5.6063   8e-04
    x     2.9974  0.1604 18.6840   0e+00
\end{verbatim}
}
\clearpage
%\newpage
\begin{figure}
\psfig{file=infl2.ps,height=9.0in}
\end{figure}
\clearpage
\begin{center}
Regression Output for Data with Leverage Point
\end{center}
{\large
\begin{verbatim}

N = 10

Residual Standard Error = 1.4046

Multiple R-Square = 0.9934

F-statistic = 1199.086 on 1 and 8 df, p-value = 0

            coef std.err  t.stat p.value 
Intercept 5.5567  0.7175  7.7444   1e-04
    x     3.0020  0.0867 34.6278   0e+00
\end{verbatim}
}
\vspace{0.5in}
\begin{center}
Regression Output for Data with LP Omitted
\end{center}
{\large
\begin{verbatim}

N = 9

Residual Standard Error = 1.5016

Multiple R-Square = 0.9717

F-statistic = 240.0134 on 1 and 7 df, p-value = 0

            coef std.err  t.stat p.value 
Intercept 5.5512  1.0909  5.0888  0.0014
    x     3.0032  0.1939 15.4924  0.0000
\end{verbatim}
}
%\newpage
\begin{figure}
\psfig{file=infl3.ps,height=9.0in}
\end{figure}
\clearpage
\begin{center}
Regression Output for Data with Influential Point
\end{center}
{\large
\begin{verbatim}

N = 10

Residual Standard Error = 6.5626

Multiple R-Square = 0.9512

F-statistic = 156.0045 on 1 and 8 df, p-value = 0

             coef std.err   t.stat p.value 
Intercept -3.8147  3.3524  -1.1379  0.2881
    x      5.0591  0.4050  12.4902  0.0000
\end{verbatim}
}
\vspace{0.5in}
\begin{center}
Regression Output for Data with IP Omitted
\end{center}
{\large
\begin{verbatim}

N = 9

Residual Standard Error = 1.5016

Multiple R-Square = 0.9717

F-statistic = 240.0134 on 1 and 7 df, p-value = 0

            coef std.err  t.stat p.value 
Intercept 5.5512  1.0909  5.0888  0.0014
    x     3.0032  0.1939 15.4924  0.0000
\end{verbatim}
}
\newpage
\begin{center}
Why Outliers Affect LS Regression
\end{center}
\vspace{0.5in}
As we've seen before, the least squares estimates are found by minimizing
the quantity Q to find the $\hat{\beta}$'s. \\
\vspace{0.35in}
\[
Q=\sum_{i=1}^{n}[y_i-f(x_{1i},x_{2i},\ldots,x_{ki};\hat{\beta}_1,\hat{\beta}_2,\ldots,\hat{\beta}_p)]^2
\] \\
\vspace{0.5in}
However, when a few deviations are much larger than the rest (in absolute value), and
all of the deviations are squared and added up, the few large deviations
tend to dominate the total. \\ 
\vspace{0.5in}
As a result, the values of the $\hat{\beta}$'s that minimize of the least squares 
criterion are often the values that eliminate large deviations from the model. 
\newpage
\begin{center}
Why Outliers Affect LS Regression
\end{center}
\vspace{0.5in}
The LSS estimates for the straight line regression function 
$y = \beta_1 + \beta_2x$ 
show in detail how both outliers and high leverage points 
affect the parameter estimates. \\ 
\vspace{0.15in}
\[ \hat{\beta}_2=\frac{\sum_{i=1}^{n}(x_i-\bar{x})(y_i-\bar{y})}{\sum_{i=1}^{n}(x_i-\bar{x})^2} \] \\
\vspace{0.15in}
\[ \hat{\beta}_1=\bar{y} - \hat{\beta}_2\bar{x} \] \\
\vspace{0.5in}
$\hat{\beta}_2$ is a linear combination of the deviations of the
data points, $y_i$, from their sample mean, $\bar{y}$, `weighted' proportional to the
deviation of $x_i$ from $\bar{x}$.  Large values of either $(x_i-\bar{x})$ or 
$(y_i-\bar{y})$ will impact the estimates of both $\beta_1$ and $\beta_2$. \\
\vspace{0.5in}
The effects of leverage points and outliers are analogous in more complicated models.
\newpage
\begin{center}
General Steps for Handling Outliers 
\end{center}
\vspace{0.5in}
The basic steps for dealing with outliers in regression analysis are:
\vspace{0.5in}
\begin{enumerate}
\item identify the outlying and influential points
\item determine the reason each of these points was observed, if possible
\item eliminate any points that can be shown to be irrelevant to the analysis
\item reanalyze the data, both with and without any remaining influential points
\end{enumerate}
\vspace{0.5in}
Steps 1 and 4 in this general procedure are statistical in nature, while steps
2 and 3 require knowledge of specific experimental details and a thorough understanding
of the underlying science.
\newpage
\begin{center}
Outliers Detection Methods
\end{center}
\vspace{0.5in}
Fortunately, outliers can usually be found using
standard model validation procedures.
However, in complicated, multivariate data, they can be 
sometimes be difficult to pick out. \\
\vspace{0.5in}
Use of more specialized outlier detection procedures 
in linear regression problems can help ensure that 
all outliers are found, and can lead to a deeper 
understanding of the data and model. \\
\vspace{0.5in}
The use and interpretation of the specialized outlier detection
methods is harder in nonlinear least squares, though logically
possible.  As a result, residual plots are the primary means of 
identifying influential points in nonlinear problems.
\newpage
\begin{center}
Two Classes of Specialized Outlier Detection Methods
\end{center}
\vspace{0.5in}
There are basically two classes of outlier detection methods: 
\vspace{0.5in}
\begin{enumerate}
\item methods that rely on the usual LS model fitting criteria,
but may (effectively) delete observations from the fit, one-at-a-time,
to help determine their influence (LS Methods), and
\vspace{0.15in}
\item methods in which the model fitting criterion is changed to
reduce the effect of outliers on model estimation (Robust Methods)
\end{enumerate}
\newpage
\begin{center}
Advantages and Disadvantages of LS Methods
\end{center}
\vspace{0.5in}
Advantages: \\
\begin{enumerate}
\item use the same software as for regular regression analysis
\item work reasonably well with small and large data sets 
\item provide a relatively large amount of information about unusual data points 
\item are part of a relatively well-developed area of statistical research
\item are computationally efficient
\end{enumerate}
\vspace{0.5in}
Disadvantages: \\
\begin{enumerate}
\item do not work well with clusters of unusual points 
\item require use of many different diagnostic statistics
\end{enumerate}
\newpage
\begin{center}
Advantages and Disadvantages of Robust Methods
\end{center}
\vspace{0.5in}
Advantages: \\
\begin{enumerate}
\item work well with arbitrary numbers and clusters of unusual points 
\item identify influential points in one step 
\end{enumerate}
\vspace{0.5in}
Disadvantages: \\
\begin{enumerate}
\item require specialized software
\item require medium-sized to large data sets
\item may not identify non-influential outlying points
\item are part of a relatively new area of statistical research
\item can be computationally difficult
\end{enumerate}
\newpage
\begin{center}
LS Methods for Outlier Detection
\end{center}
\vspace{0.5in}
Four of the main LS statistics for outlier detection are:
\begin{description}
\item[leverages] - which are used to identify high
leverage points (unusual predictor variable values), regardless of their
influence,
\item[studentized deleted residuals] -  which are used
to identify outliers (unusual response variable values)
which may or may not be influential,
\item[Cook's distances] - which are used to identify points that
influence one or more of the estimated parameters, and
\item[DFFITS] - which are used to identify points which influence
one or more of the predicted values. 
\end{description}
\newpage
\begin{center}
Leverage Values (aka Hat Values)
\end{center}
\vspace{0.5in}
Leverage values are the diagonal values of the Hat matrix, $H$,
which is given by:
\[ H = X(X^TX)^{-1}X^T \] \\
where $X$ is the matrix with $n$ rows and $p = k+1$ columns,
the first of which is a column of $1$'s for the intercept term, and 
each of the other columns given by the values of 
one of the $k$ variables in the model, 
 \[ X = [1 | x_1 | x_2 | \ldots | x_k] \] \\
\vspace{0.5in}
For the straight line model 
\[ h_{ii} = \frac{1}{n} + \frac{(x_i-\bar{x})^2}{\sum_{i=1}^n(x_i-\bar{x})^2} \] \\
\vspace{0.5in}
Leverages for more complicated models are conceptually analogous to the straight line case. 
\newpage
\begin{center}
Leverage Values (aka Hat Values)
\end{center}
\vspace{0.5in}
The sum of the leverage values is $p$, the number of parameters in the model, for all multivariable linear regression functions. \\
\vspace{0.5in}
Therefore, when an individual leverage value is much larger than the mean leverage $p/n$, it
indicates a high leverage point. \\
\vspace{0.5in}
A good cut-off for leverage values is $2.5p/n$.  Points with values larger than this are
likely to be contributing to the fitted model more heavily than most. \\
\vspace{0.5in}
High leverage doesn't necessarily indicate an influential point, however.  A high leverage
point with a $y$ value that is right on target will produce a fit that is similar to 
one using only low leverage points.
\newpage
\begin{center}
Studentized Deleted Residuals
\end{center}
\vspace{0.5in}
The $i^{th}$ studentized deleted residual is conceptually obtained by: 
\begin{enumerate}
\item deleting the $i^{th}$ point from the data,
\item fitting the model with the other $n-1$ observations,
\item computing the deleted residual, which is the difference of the
$i^{th}$ response and corresponding prediction from the fit made without the $i^{th}$
response.
\item computing the residual standard deviation using the
fit to the $n-1$ other points, and
\item dividing the deleted residual by its standard deviation, which
depends on the residual standard deviation with the $i^{th}$ point
deleted and the $i^{th}$ leverage value computed using all of the data. \\
\end{enumerate}
\newpage
\begin{center}
Studentized Deleted Residuals
\end{center}
\vspace{0.5in}
The conceptual formula for the $i^{th}$ studentized deleted residual is 
\vspace{0.15in}
\[  T_i = \frac{y_i-\hat{y}_{(-i)}(x_{1i},\ldots,x_{ki})}{s_{(-i)}/\sqrt{1-h_{ii}}} \] \\
\vspace{0.15in}
where $\hat{y}_{(-i)}(x_{1i},\ldots,x_{ki})$ is the predicted value based on the fit without
the $i^{th}$ point, and $s_{(-i)}$ is the corresponding residual standard deviation.  \\
\vspace{0.5in}
The computational formula is:
\vspace{0.15in}
\[ T_i = e_i\sqrt{\frac{n-p-1}{(n-p)s^2(1-h_{ii}) - e_i^2}} \] \\
\vspace{0.15in}
where the computations are done using all of the data.
\newpage
\begin{center}
Studentized Deleted Residuals
\end{center}
\vspace{0.5in}
Technically, $T_i$ has a t distribution with $n-p-1$ degrees of 
freedom, and t distribution cut-off values can be used to determine
if a particular point is an outlier. \\
\vspace{0.5in}
Practically speaking, 
studentized deleted residuals greater than $2.5$ in absolute value 
indicate that a point is likely to be an outlier. \\
\vspace{0.5in}
Remember, however, that large studentized deleted residuals do not 
necessarily offer any information about the influence
of the outlier on the fit of the model.
\newpage
\begin{center}
Cook's Distance 
\end{center}
\vspace{0.5in}
Conceptually, the $i^{th}$ Cook's Distance measures the difference
between the values of the parameters estimated using all $n$ observations
and the values of the parameters estimated without the $i^{th}$ observation. \\
\vspace{0.5in}
It provides an aggregate measure of the difference for all of the parameters
simultaneously. \\
\vspace{0.5in}
The computational formula for Cook's Distance is: 
\vspace{0.15in}
\[ c_i = \frac{h_{ii}}{p}(\frac{e_i}{s(1-h_{ii})})^2 \] \\
\vspace{0.15in}
computed using all of the data.
\newpage
\begin{center}
Cook's Distance 
\end{center}
\vspace{0.5in}
Points with Cook's Distance values greater than the 50\% F distribution cut-off
with $p$ and $n-p$ degrees of freedom are likely to influence the estimated
parameters in the model. \\
\vspace{0.5in}
That is, if a point with a large Cook's Distance is eliminated from the 
data set, the values of one or more of the parameters will change 
significantly. \\
\vspace{0.5in}
Influential points should be studied carefully
to determine if they are valid data points.  If not, they should be excluded from the
analysis.  If there is no evidence that the influential points are invalid, further
work will be required to obtain conclusive results from the data.
\newpage
\begin{center}
DFFITS
\end{center}
\vspace{0.5in}
Conceptually, the $i^{th}$ DFFITS value measures the difference
between the $i^{th}$ predicted value estimated using all $n$ observations
and estimated without the $i^{th}$ observation. \\
\vspace{0.5in}
The conceptual formula for DFFITS is: 
\vspace{0.15in}
\[ \mbox{DFFITS}_i = \frac{\hat{y}_i(x_{1i},\ldots,x_{ki}) - \hat{y}_{(-i)}(x_{1i},\ldots,x_{ki})}{s_{(-i)}\sqrt{h_{ii}}}  \] \\
\vspace{0.5in}
The computational formula for DFFITS is: 
\vspace{0.15in}
\[ \mbox{DFFITS}_i = T_i\sqrt{\frac{h_{ii}}{1-h_{ii}}} \] \\
\vspace{0.15in}
computed using all of the data.
\newpage
\begin{center}
DFFITS
\end{center}
\vspace{0.5in}
Points with DFFITS values greater than $2.5\sqrt{p/n}$ are likely to influence 
the predicted values from the model. \\
\vspace{0.5in}
That is, if a point with a large DFFITS value is eliminated from the data set, the 
value of the one or more of the predicted values will change significantly. \\
\vspace{0.5in}
As with influential points identified using Cook's Distance, the validity of 
points picked out by DFFITS should be examined carefully.  If assignable causes
are found that indicate the point is not valid, it should be dropped from the
analysis.  If not, further work will required to obtain conclusive results.
%newpage
\begin{figure}
\psfig{file=infl4.ps,height=9.0in}
\end{figure}
%newpage
\begin{figure}
\psfig{file=infl5.ps,height=9.0in}
\end{figure}
%newpage
\begin{figure}
\psfig{file=infl6.ps,width=6.5in}
\end{figure}
%newpage
\begin{figure}
\psfig{file=infl7.ps,height=9.0in}
\end{figure}
%newpage
\begin{figure}
\psfig{file=infl8.ps,height=9.0in}
\end{figure}
\clearpage
\newpage
\begin{center}
Regression Output Using n Points
\end{center}
{\large
\begin{verbatim}

N = 30

Residual Standard Error = 0.1859

Multiple R-Square = 1

F-statistic = 1121069 on 2 and 27 df, p-value = 0

              coef std.err    t.stat p.value 
Intercept   4.7251  0.2153   21.9480       0
   x1     -11.9502  0.0345 -346.3767       0
   x2      10.0313  0.0674  148.9384       0
\end{verbatim}
}
\vspace{0.5in}
\begin{center}
Regression Output Using with the High Leverage Point Deleted
\end{center}
{\large
\begin{verbatim}

N = 29

Residual Standard Error = 0.1895

Multiple R-Square = 1

F-statistic = 1065271 on 2 and 26 df, p-value = 0

              coef std.err    t.stat p.value 
Intercept   4.6989  0.3511   13.3843       0
   x1     -11.9458  0.0577 -206.9284       0
   x2      10.0405  0.1181   85.0499       0
\end{verbatim}
}
%newpage
\begin{figure}
\psfig{file=infl9.ps,height=9.0in}
\end{figure}
%newpage
\begin{figure}
\psfig{file=infl10.ps,height=9.0in}
\end{figure}
%newpage
\begin{figure}
\psfig{file=infl11.ps,height=9.0in}
\end{figure}
%newpage
\begin{figure}
\psfig{file=infl12.ps,width=6.5in}
\end{figure}
%newpage
\begin{figure}
\psfig{file=infl13.ps,height=9.0in}
\end{figure}
%newpage
\begin{figure}
\psfig{file=infl14.ps,height=9.0in}
\end{figure}
\clearpage
\newpage
\begin{center}
Regression Output Using n Points
\end{center}
{\large
\begin{verbatim}

N = 30

Residual Standard Error = 1.2939

Multiple R-Square = 0.9994

F-statistic = 22959.24 on 2 and 27 df, p-value = 0

              coef std.err   t.stat p.value 
Intercept  14.3252  1.4980   9.5630       0
   x1     -13.5540  0.2401 -56.4605       0
   x2       6.6670  0.4686  14.2261       0
\end{verbatim}
}
\vspace{0.5in}
\begin{center}
Regression Output Using with the Influential Point Deleted
\end{center}
{\large
\begin{verbatim}

N = 29

Residual Standard Error = 0.1895

Multiple R-Square = 1

F-statistic = 1065271 on 2 and 26 df, p-value = 0

              coef std.err    t.stat p.value 
Intercept   4.6989  0.3511   13.3843       0
   x1     -11.9458  0.0577 -206.9284       0
   x2      10.0405  0.1181   85.0499       0
\end{verbatim}
}
\clearpage
%newpage
\begin{figure}
\psfig{file=infl15.ps,width=6.5in}
\end{figure}
%newpage
\begin{figure}
\psfig{file=chwirut_out1.ps,width=6.5in}
\end{figure}
%newpage
\begin{figure}
\psfig{file=chwirut_out2.ps,width=6.5in}
\end{figure}
%newpage
\begin{figure}
\psfig{file=chwirut_out3.ps,width=6.5in}
\end{figure}
\clearpage
\newpage
\begin{center}
Ultrasonic Data Regression Output: \\ n Points
\end{center}
{\large
\begin{verbatim}

Formula: sqrt(ur) ~ exp( - b1 * md)/(b2 + b3 * md)

Parameters:

        Value Std. Error   t value 
b1 -0.0154274 0.00861101  -1.79159
b2  0.0806725 0.00150574  53.57670
b3  0.0638570 0.00288001  22.17250

Residual standard error: 0.29715 on 211 degrees of freedom
\end{verbatim}
}
\vspace{0.5in}
\begin{center}
Ultrasonic Data Regression Output: \\ n-2 Points
\end{center}
{\large
\begin{verbatim}

Formula: sqrt(ur) ~ exp( - b1 * md)/(b2 + b3 * md)

Parameters:

        Value Std. Error   t value 
b1 -0.0155610 0.00799912  -1.94533
b2  0.0803004 0.00140502  57.15260
b3  0.0644030 0.00268723  23.96630

Residual standard error: 0.276145 on 209 degrees of freedom
\end{verbatim}
}
\clearpage
%newpage
\begin{figure}
\psfig{file=chwirut_out4.ps,width=6.5in}
\end{figure}
\clearpage
\newpage
\markboth{}{}
\begin{center}
Section 2: Summary
\end{center}
\vspace{0.5in}
Two problems that often occur when carrying out
a regression analysis are:
\vspace{0.15in}
\begin{enumerate}
\item finding non-constant standard deviations 
across different predictor variable values, and
\vspace{0.15in}
\item finding outliers, high leverage points, and/or influential
points in the data.
\end{enumerate}
\vspace{0.5in}
Outliers, leverage points and influential points can be identified 
using either graphical residual analysis or more specialized detection 
methods. \\
\vspace{0.5in}
Non-constant standard deviation across the predictors can be 
corrected using either transformation of the data or weighted least
squares.
\vspace{0.5in}

%close out general text format for slides
}
}
% end close out
\end{document}
